%=======================================================================
%  Q6_DEFAULT_Rev-.tex
%  Paper 6 of the R(x)C(x)H(x)O series
%  Landau-Gauge W+/- Propagator from the b-Channel:
%  Gauge Fixing, Gribov Copies, and the Cross-Projected
%  Correlator in SU(2)
%
%  Draft assembled from section drafts s2-s7, the Section 8 algebra
%  output, and the Paper 6 outline. Sections 1, 8, 9, 10 and the
%  appendices are drafted from the outline and editor briefs.
%  Status: working draft for internal review.
%=======================================================================
\documentclass[11pt,a4paper]{article}

\usepackage[T1]{fontenc}
\usepackage[utf8]{inputenc}
\IfFileExists{lmodern.sty}{\usepackage{lmodern}}{}
\usepackage{amsmath,amssymb,amsthm}
\usepackage{mathtools}
\usepackage{bm}
\usepackage{geometry}
\geometry{margin=1in}
\usepackage{booktabs}
\usepackage{array}
\usepackage{enumitem}
\usepackage{graphicx}
\usepackage{xcolor}
\usepackage[hidelinks]{hyperref}
\usepackage{cleveref}

%-----------------------------------------------------------------------
%  Theorem environments
%-----------------------------------------------------------------------
\theoremstyle{plain}
\newtheorem{theorem}{Theorem}
\newtheorem{lemma}{Lemma}
\newtheorem{proposition}{Proposition}
\newtheorem{corollary}{Corollary}

\theoremstyle{definition}
\newtheorem{definition}{Definition}

\theoremstyle{remark}
\newtheorem*{remark}{Remark}
\newtheorem*{note}{Note}

%-----------------------------------------------------------------------
%  Notation macros (recurring objects of the series)
%-----------------------------------------------------------------------
\newcommand{\Cl}{\mathrm{Cl}}
\newcommand{\SU}{\mathrm{SU}}
\newcommand{\SO}{\mathrm{SO}}
\newcommand{\Tr}{\operatorname{Tr}}
\newcommand{\Ree}{\operatorname{Re}}
\newcommand{\Imm}{\operatorname{Im}}
\renewcommand{\Re}{\operatorname{Re}}
\newcommand{\Gcross}{G_{\mathrm{cross}}}
\newcommand{\Gcrossgf}{G_{\mathrm{cross}}^{\mathrm{gf}}}
\newcommand{\Gdiag}{G_{\mathrm{diag}}}
\newcommand{\Gdiaggf}{G_{\mathrm{diag}}^{\mathrm{gf}}}
\newcommand{\Gfull}{G_{\mathrm{full}}}
\newcommand{\DW}{D_{W^{\pm}}}
\newcommand{\DWt}{\widetilde{D}_{W^{\pm}}}
\newcommand{\Scross}{S_{\mathrm{cross}}}
\newcommand{\be}{\beta_{\mathrm{eff}}}
\newcommand{\bc}{\beta_{c}}
\newcommand{\epsg}{\varepsilon_{\mathrm{gauge}}}
\newcommand{\FL}{F_{L}}
\newcommand{\Ncfg}{N_{\mathrm{cfg}}}
\newcommand{\Nt}{N_{t}}
\newcommand{\Pew}{P^{\mathrm{EW}}}
\newcommand{\half}{\tfrac{1}{2}}
\newcommand{\Id}{\mathbb{I}}
\newcommand{\PL}{P_{L}}
\newcommand{\PR}{P_{R}}
\newcommand{\gfive}{\gamma_{5}}
\newcommand{\Lcond}{\partial_{\mu}A_{\mu}}

%-----------------------------------------------------------------------
\title{\bfseries Landau-Gauge $W^{\pm}$ Propagator from the
$b$-Channel:\\[2pt]
Gauge Fixing, Gribov Copies, and the Cross-Projected\\
Correlator in $\SU(2)$}

\author{T. Firestone}

\date{}

%=======================================================================
\begin{document}
\maketitle

\begin{abstract}
\noindent
Paper 5 established the exact algebraic decomposition
$\Gdiag(t)+\Gcross(t)=\Gfull(t)$ in the $\Cl(0,2)\cong\mathbb{H}$
framework and identified $\Gcross(t)=8\sum_{x,\mu}\Re[b_t\bar b_0]$
as the $b$-channel bilinear carrying the charged ($W^{\pm}$) content
of the embedded $\SU(2)$ link. That identification was, however,
gauge dependent: $\Gcross$ is not gauge invariant, no gauge
prescription had been imposed, and no propagator had been extracted.
This paper closes the gap. We implement Landau gauge fixing on the
lattice through steepest-descent and Fourier-accelerated algorithms,
prove a convergence-rate lemma (Lemma 4), and characterise the Gribov
copy systematic through a multi-start best-copy strategy and an
explicit variation bound (Proposition 7). The central result,
Theorem 2, promotes the formal perturbative identification of
Paper 5 to a fully specified, error-budgeted quantity: the gauge-fixed correlator
$\Gcrossgf(t)$ yields the zero-spatial-momentum $W^{\pm}$ temporal
propagator $\DW(t)$ with a complete three-term error budget. Corollary 6
characterises the sign structure of the ensemble-averaged correlator
and defines a coupling threshold $\bc$ for the onset of
non-perturbative sign violation, which is studied empirically through
an extended ensemble programme spanning $N\in\{6,8,12,16\}$ and
$\be\in\{2.0,\ldots,20.0\}$ with a $1/N^{2}$ finite-volume
extrapolation. We establish $\SU(2)$ colour symmetry between the
$b$-channel propagator and the full Landau-gauge gluon propagator,
enabling comparison to the $\SU(2)$ reference dataset of
Cucchieri--Mendes (with Bogolubsky et al.\ providing a qualitative
$\SU(3)$ confirmation of the same decoupling behaviour). Finally, the fermionic sector is set up: we define the
cross-projected propagator $\Scross=\Pew_{+}S\,\Pew_{-}$, prove that it
vanishes identically for an $\SU(2)$ singlet (Lemma 5), and specify the
weak-coupling verification (checks D28/D29) that the doublet
cross-projected propagator is non-zero. Throughout, the embedding does
not by itself distinguish $\SU(2)_{L}$ from $\SU(2)_{R}$; this requires
external fermionic input, the full non-perturbative test of which is
deferred to Paper 7.
\end{abstract}

%=======================================================================
\section{Introduction}
\label{sec:intro}

Paper 5 of this series established, within the
$\Cl(0,2)\cong\mathbb{H}$ framework, the exact identity
\begin{equation}
  \Gdiag(t)+\Gcross(t)=\Gfull(t),
  \label{eq:thm1}
\end{equation}
(Theorem 1, Paper 5) together with the Two-Observable Decomposition
(Corollary 5, Paper 5), in which the cross-channel correlator
\begin{equation}
  \Gcross(t)=8\sum_{x,\mu}\Re\!\bigl[b_t(x,\mu)\,\bar b_0(x,\mu)\bigr]
  \label{eq:gcrossdef}
\end{equation}
is identified as the $b$-channel bilinear, with $b=(U)_{01}$ the
off-diagonal element of the $\SU(2)$ link
$U=\bigl(\begin{smallmatrix} a & b\\ -\bar b & \bar a
\end{smallmatrix}\bigr)$, $|a|^2+|b|^2=1$. At weak coupling
$b=(ig/\sqrt2)W^{+}_{\mu}+O(g^{3})$, so $\Gcross$ carries the charged
($W^{\pm}$) content of the embedded gauge field, while the diagonal
correlator $\Gdiag$ carries the neutral $W^{3}/A^{3}$ content
(Paper 5, \S7.3).

\paragraph{The inherited gap.}
The Paper 5 identification was incomplete in six specific respects.
(i)~$\Gcross$ is not gauge invariant (gauge non-invariance, Paper~5 \S8.1), yet no
gauge prescription had been fixed. (ii)~The formal identification of
Paper 5 \S8.3 applied to spatial links only, because Papers 4--5 worked
in temporal gauge $U_0(x)=\Id_2$. (iii)~No control of Gribov copy
ambiguity had been established. (iv)~The $W^{\pm}$ temporal propagator
$\DW(t)$ had not been extracted from any ensemble. (v)~No error budget
existed. (vi)~No comparison to independent lattice data had been made.
Paper 6 addresses all six, organising its verification checks D19--D29
against them.

\paragraph{What this paper establishes.}
The centrepiece is Theorem 2 (\Cref{sec:theorem2}), which promotes the
Paper 5 \S8.3 formal identification to a fully specified, error-budgeted
quantity. The route requires: a Landau gauge functional summed over all
four link directions, lifting the temporal-link restriction of Papers
4--5 (\Cref{sec:gaugefix}); a convergence-rate result for the
gauge-fixing functional (Lemma 4); a multi-start best-copy strategy
with an explicit Gribov variation bound (Proposition 7,
\Cref{sec:gribov}); and an extended ensemble programme with
finite-volume extrapolation (\Cref{sec:ensembles}). Corollary 6
characterises the sign structure of $\Gcrossgf$, and
\Cref{sec:signstructure} delivers its empirical content. Colour
symmetry between the $b$-channel and the full gluon propagator
(\Cref{sec:comparison}) enables comparison to the $\SU(2)$ reference
dataset of Cucchieri--Mendes~\cite{CM} (Bogolubsky et al.~\cite{Bo} confirm
the same decoupling in $\SU(3)$). The
fermionic groundwork (\Cref{sec:fermionic}) defines the cross-projected
propagator and proves the singlet vanishing (Lemma 5).

\paragraph{Scope limits.}
The following hard scope limits are carried forward from Papers 4--5
and remain in force except where explicitly addressed: (SL1) no
$\SU(2)_L/\SU(2)_R$ distinction without fermion content; (SL2) no
chiral fermion claims without a Nielsen--Ninomiya-bypass action; (SL3)
no $W/Z$ mass ratio, Weinberg angle, or $U(1)_Y$ mixing; (SL4) no Higgs
mechanism; (SL5) no non-perturbative $\Gcross\leftrightarrow W^{\pm}$
identification without the Landau-gauge apparatus. SL5 is discharged in
\Cref{sec:gribov}; SL1 and SL2 are addressed (not fully retired) in
\Cref{sec:fermionic}.

\begin{quote}
\itshape
L/R caveat (verbatim, first occurrence): The embedding does not, by
itself, distinguish $\SU(2)_L$ from $\SU(2)_R$; this requires external
fermionic input.
\end{quote}

\paragraph{Infrastructure.}
The simulation functions implemented in this paper are, in paper-local
numbering, P6-F1 (steepest descent) and P6-F2 (Fourier acceleration) in
\Cref{sec:gaugefix}; P6-F3 (best-copy strategy) in \Cref{sec:gribov};
P6-F4 (gauge-fixed correlator) in \Cref{sec:theorem2}; P6-F6 through
P6-F9 (sign analysis) in \Cref{sec:signstructure}; and P6-F10 through
P6-F13 (fermionic sector) in \Cref{sec:fermionic}. New numbered results
are Lemma 4 (\Cref{sec:gaugefix}), Proposition 7 (\Cref{sec:gribov}),
Theorem 2 and Corollary 6 (\Cref{sec:theorem2}), and Lemma 5
(\Cref{sec:fermionic}).

%=======================================================================
\section{Landau Gauge Fixing: Algorithm and Implementation}
\label{sec:gaugefix}

The algebraic programme of Papers 4--5 established $\Gcross$ exactly as
a $b$-channel observable but left it gauge dependent (gauge non-invariance,
Paper~5 \S8.1). Before a propagator can be extracted, a gauge prescription
must be fixed. This section develops the Landau gauge-fixing machinery:
the functional $\FL$ is defined (\S\ref{ssec:functional}), the lattice
Landau condition and convergence criterion are derived
(\S\ref{ssec:landaucond}), the steepest-descent algorithm is specified
and its monotone-increase property proved (\S\ref{ssec:SD}), the
Fourier-accelerated algorithm is specified (\S\ref{ssec:FA}), and
Lemma 4 establishes convergence rates and the regime of preference
(\S\ref{ssec:lemma4}). No measurement of $\Gcrossgf$ is performed here;
that is \Cref{sec:theorem2}. The multiplicity of local maxima of $\FL$
(Gribov copies) is treated in \Cref{sec:gribov}.

\subsection{The Landau Gauge Functional}
\label{ssec:functional}

\begin{definition}[Landau Gauge Functional]
For a set of $\SU(2)$ gauge transformations $\{g(x)\}$ on the lattice
$\Lambda=\{0,\ldots,N-1\}^3\times\{0,\ldots,\Nt-1\}$ with periodic
boundary conditions, the Landau gauge functional is
\begin{equation}
  \FL[\{g(x)\}]=\Re\sum_{x,\mu}\Tr\bigl[U^{g}_{\mu}(x)\bigr],
  \label{eq:FLdef}
\end{equation}
where $U^{g}_{\mu}(x)=g(x)\,U_{\mu}(x)\,g^{\dagger}(x+\hat\mu)$ and the
sum runs over all links: $x\in\Lambda$ and $\mu\in\{0,1,2,3\}$,
including the temporal direction $\mu=0$.
\end{definition}

Since $\Tr[U^g_\mu]+\Tr[(U^g_\mu)^\dagger]=2\Re\Tr[U^g_\mu]$, the
functional is real. Each link contributes at most $\Re\Tr[\Id_2]=2$ and
at least $-2$, giving the uniform bounds
\begin{equation}
  0\le \FL[\{g(x)\}]\le \FL^{\max}:=2\,dV,
  \label{eq:FLbounds}
\end{equation}
with $d=4$ and $V=N^4$, so $\FL^{\max}=8V$.

\begin{remark}
Papers 4--5 imposed temporal gauge $U_0(x)=\Id_2$ before measurement,
restricting $\FL$ to spatial links (maximum $6V$). Paper 5 \S8.3 noted
this as a caveat. Paper 6 lifts the restriction: the temporal links are
live variables and \eqref{eq:FLdef} sums over all four directions,
matching the standard prescription of \cite{CM,Bo} and enabling the
quantitative comparison of \Cref{sec:comparison}.
\end{remark}

Landau gauge is the choice of representatives $\{g^{*}(x)\}$ globally
maximising $\FL$,
\begin{equation}
  \{g^{*}(x)\}=\operatorname*{argmax}_{\{g(x)\}\in\SU(2)^{V}}
  \FL[\{g(x)\}].
  \label{eq:argmax}
\end{equation}
Global maximisation is generically obstructed by multiple local maxima
on each gauge orbit --- the Gribov copies~\cite{Gri,Zwa} --- treated in
\Cref{sec:gribov}.

\subsection{The Lattice Landau Condition}
\label{ssec:landaucond}

\begin{definition}[Lattice Gauge Field]
For a gauge-transformed link, the lattice gauge field is
\begin{equation}
  A^{a}_{\mu}(x)=\frac{1}{2i}\Tr\!\bigl[T^{a}
  \bigl(U^{g}_{\mu}(x)-(U^{g}_{\mu}(x))^{\dagger}\bigr)\bigr],
  \label{eq:Afield}
\end{equation}
with $T^{a}=\sigma^{a}/2$ the $\SU(2)$ generators.
\end{definition}

\begin{definition}[Lattice Divergence and Convergence Measure]
The lattice divergence and the global convergence measure are
\begin{align}
  (\Lcond)^{a}(x)&=\sum_{\mu=0}^{3}
  \bigl[A^{a}_{\mu}(x)-A^{a}_{\mu}(x-\hat\mu)\bigr],
  \label{eq:divergence}\\
  \Theta&=\frac{1}{d\,N_c\,V}\sum_{x}\sum_{a=1}^{3}
  \bigl[(\Lcond)^{a}(x)\bigr]^{2},
  \label{eq:Theta}
\end{align}
with $d=4$, $N_c=2$. The factor $d=4$ in the denominator is essential;
using $d=3$ would inflate $\Theta$ by $33\%$ and cause premature
termination.
\end{definition}

The lattice Landau condition is $\Theta=0$; the numerical criterion
adopted throughout is
\begin{equation}
  \Theta<\epsg=10^{-14}.
  \label{eq:epsgauge}
\end{equation}
Under $g(x)\to(\Id+i\varepsilon\,\omega^{a}(x)T^{a})g(x)$,
\begin{equation}
  \delta\FL=-2\varepsilon\sum_{x,a}\omega^{a}(x)\,(\Lcond)^{a}(x)
  +O(\varepsilon^{2}),
  \label{eq:firstvar}
\end{equation}
so stationarity requires $(\Lcond)^{a}(x)=0$ for all $x,a$ --- the
lattice Landau condition.

\subsection{Steepest-Descent Gauge Fixing (P6-F1)}
\label{ssec:SD}

\begin{definition}[Local Staple]
With all $g(y)$, $y\ne x$, fixed, the site-$x$ contribution to $\FL$ is
$f_x[g(x)]=\Re\Tr[g(x)W(x)]$ with staple
\begin{equation}
  W(x)=\sum_{\mu=0}^{3}\bigl[U^{g}_{\mu}(x)+(U^{g}_{\mu}(x-\hat\mu))^{\dagger}\bigr],
  \label{eq:staple}
\end{equation}
parametrised as $W(x)=w_0\Id_2+i\,\bm{w}\cdot\bm{\sigma}$ with
$|W(x)|=(w_0^2+|\bm{w}|^2)^{1/2}$.
\end{definition}

\begin{proposition}[Local Maximisation on $\SU(2)$]
For fixed $W(x)$ with $|W(x)|>0$, the unique maximiser of
$\Re\Tr[g(x)W(x)]$ over $g(x)\in\SU(2)$ is
$g^{*}(x)=W^{\dagger}(x)/|W(x)|$.
\end{proposition}

\begin{proof}
Write $g=g_0\Id_2+i\bm{g}\cdot\bm{\sigma}$, $g_0^2+|\bm{g}|^2=1$. Using
$\Tr[\Id_2]=2$ and $\Tr[\sigma^a\sigma^b]=2\delta^{ab}$,
\[
  \Re\Tr[g(x)W(x)]=2(g_0w_0-\bm{g}\cdot\bm{w})
  =2\,(g_0,\bm{g})\cdot(w_0,-\bm{w}).
\]
By Cauchy--Schwarz with $|(g_0,\bm{g})|=1$, this is at most $2|W(x)|$,
with equality iff $(g_0,\bm{g})\propto(w_0,-\bm{w})$, i.e.
$g^{*}(x)=W^{\dagger}(x)/|W(x)|$; $\det g^*=1$ so $g^*\in\SU(2)$.
\end{proof}

\paragraph{Algorithm P6-F1 (steepest descent).}
Initialise $g^{(0)}(x)=R(x)$, Haar-random $\SU(2)$ per site. Iterate:
for each site $x$ compute $W^{(k)}(x)$ and set
$g^{(k+1)}(x)=(W^{(k)}(x))^{\dagger}/|W^{(k)}(x)|$, updating the
adjacent links; after each full sweep compute $\Theta$; stop when
$\Theta<\epsg$, recording the sweep count $k^{*}$; flag non-convergence
beyond $k_{\max}=10\,000$. Monotone increase of $\FL$ is Lemma 4(i). An
over-relaxation variant ($\omega\approx1.7$--$1.9$) reduces the sweep
count by $O(N)$ but is not used for the $N=6$ ensembles.

\subsection{Fourier-Accelerated Gauge Fixing (P6-F2)}
\label{ssec:FA}

Steepest descent converges slowly for long-wavelength modes
($\rho_{\mathrm{SD}}\approx1-\pi^2/N^2$). The Fourier-accelerated (FA)
algorithm preconditions the update in momentum space, reducing the
sweep count to $O(N)$ while preserving monotone increase in the
linearised regime; it is standard for large volumes~\cite{CM}.

\paragraph{Algorithm P6-F2.}
Each pass: compute $A^a_\mu(x)$, the divergence \eqref{eq:divergence},
and $\Theta$; stop if $\Theta<\epsg$. Otherwise Fourier transform the
divergence,
\begin{equation}
  \tilde p^{a}(p)=\sum_{x}e^{2\pi i\,p\cdot x/N}(\Lcond)^{a}(x),\qquad
  \tilde p^{a}(0):=0,
  \label{eq:fft}
\end{equation}
precondition with the lattice Laplacian eigenvalue
$f(p)=\sum_\mu 4\sin^2(\pi p_\mu/N)$,
\begin{equation}
  \delta\tilde\omega^{a}(p)=\tilde p^{a}(p)/f(p),
  \label{eq:precond}
\end{equation}
invert the transform, exponentiate
$h(x)=\exp(i\alpha\,\delta\omega^a(x)T^a)$ with step size $\alpha$, and
update $g^{(k+1)}(x)=h(x)g^{(k)}(x)$. The step size is tuned per
$\be$ ($\alpha=0.08$ at $\be\ge9$ down to $\alpha=0.20$ at
$\be\le3$), with optimal value $\alpha_{\mathrm{opt}}=2/(f_{\min}+f_{\max})$,
$f_{\max}=16$. The zero-mode projection removes the global rotation.
\texttt{scipy.fft} (1.16.3) implements the transforms.

\subsection{Lemma 4: Convergence Rate of \texorpdfstring{$\FL$}{F\_L}}
\label{ssec:lemma4}

\setcounter{lemma}{3}
\begin{lemma}[Convergence Rate of $\FL$ on $\SU(2)$ Ensembles]
Let $\Lambda$ be hypercubic with $N$ sites per direction, $\Nt=N$,
$V=N^4$; $\{U_\mu(x)\}$ a thermalised $\SU(2)$ ensemble at coupling
$\be$; $\epsg=10^{-14}$.
\begin{enumerate}[label=(\roman*)]
\item \textbf{Monotone increase.} Every site update in P6-F1 satisfies
$\FL[\{g^{(k+1)}\}]\ge\FL[\{g^{(k)}\}]$, with equality iff $\Theta=0$
everywhere. The algorithm converges in finite $k^{*}$ to a local
maximum of $\FL$.
\item \textbf{$\be$ dependence.}
$N_{\mathrm{iter}}^{\mathrm{SD}}\approx(N^2/\pi^2)\,
(1+O(8/\be))\,|\log\epsg|=O\!\bigl(V^{1/2}(1+c/\be)|\log\epsg|\bigr)$.
\item \textbf{Fourier acceleration.}
$N_{\mathrm{iter}}^{\mathrm{FA}}=O\!\bigl(V^{1/4}(1+c/\be)|\log\epsg|\bigr)$
passes, each of cost $O(V\log V)$. The speedup over SD is $O(N/\log N)$,
exceeding unity for $N\ge12$.
\end{enumerate}
\end{lemma}

\begin{proof}
\textit{(i)} Fixing $g(y)$, $y\ne x$, the update sets
$\Re\Tr[g^{(k+1)}(x)W(x)]=2|W(x)|$, the global maximum, while
$\Re\Tr[g^{(k)}(x)W(x)]\le 2|W(x)|$; hence $\Delta F_x\ge0$ at every
site and $\FL^{(k+1)}\ge\FL^{(k)}$. Equality requires $W(x)\propto
g^{(k)}(x)^{\dagger}$ for all $x$, the stationarity condition
$\Theta=0$. Bounded above by $8V$ and non-decreasing, the sequence
converges to a local maximum.

\textit{(ii)} Near a local maximum, the second variation of $\FL$ is
governed by the Faddeev--Popov operator $M_{\mathrm{FP}}$, which in the
$\be\to\infty$ limit reduces to the lattice Laplacian with eigenvalues
$\lambda_p=\sum_\mu 4\sin^2(\pi p_\mu/N)$, smallest non-zero
$\lambda_1\approx4\pi^2/N^2$, largest $\lambda_{\max}=16$, condition
number $\kappa\approx(4/\pi^2)N^2$. The SD reduction factor per sweep is
$\rho_{\mathrm{SD}}\approx1-\pi^2/N^2$, giving
$N_{\mathrm{iter}}^{\mathrm{SD}}\approx(N^2/\pi^2)|\log(\epsg/\Theta_0)|$,
with the $O(g^2)$ correction at finite $\be$ multiplying by
$(1+O(8/\be))$. At strong coupling ($\be\lesssim3$) near-zero
Faddeev--Popov modes signal Gribov-horizon proximity; this is a physical
effect, flagged in D19.

\textit{(iii)} FA inverts the Laplacian part of $M_{\mathrm{FP}}$ each
pass, so the spectral \emph{gap} square-roots,
$(1-\rho_{\mathrm{FA}})\approx(1-\rho_{\mathrm{SD}})^{1/2}$; with
$1-\rho_{\mathrm{SD}}\approx\pi^2/N^2$ this gives
$1-\rho_{\mathrm{FA}}\approx\pi/N$ and $O(N)$ passes. SD costs
$O(V^{3/2})$ total operations, FA costs $O(V^{5/4}\log V)$; the speedup
is $O(V^{1/4}/\log V)=O(N/\log N)$, giving heuristically (hidden constants
set to unity) $S(12)\approx1.20$, $S(16)\approx1.44$.
\end{proof}

\paragraph{Checks D19--D21.}
D19 verifies $\Theta<10^{-14}$ after every gauge-fixing call (expected
pass rate $\ge95\%$ at $\be\in\{5,\ldots,12\}$; lower at strong coupling
with stalling flagged, not failed). D20 verifies the per-site condition
$\max_{x,a}|(\Lcond)^a(x)|<\varepsilon_{\mathrm{site}}=\sqrt{8V\cdot10^{-14}}$.
D21 verifies (a)~$\FL^{\mathrm{gf}}>\FL^{\mathrm{initial}}$,
(b)~$\FL^{\mathrm{gf}}/(8V)\in[R_{\min}(\be),1]$, and
(c)~extensive volume scaling within $5\%$.

\subsection{Algorithm Comparison and Regime Decision}
\label{ssec:regime}

The re-processed Paper 4/5 ensembles ($N=6$, $\be\in\{5,6,9,12\}$) are
gauge-fixed with P6-F1; the marginal FA speedup at $N=6$ does not justify
the FFT infrastructure. All extended ensembles ($N\in\{8,12,16\}$,
$\be\in\{2,\ldots,20\}$) use P6-F2. In both cases $\Theta<10^{-14}$ with
$d=4$ in \eqref{eq:Theta}, and the best-copy strategy P6-F3
(\Cref{sec:gribov}) with $\ge5$ random starts precedes any measurement.

%=======================================================================
\section{Gribov Copy Analysis and Systematic Uncertainty}
\label{sec:gribov}

A local maximum of $\FL$ need not be unique on a gauge orbit. The
existence of multiple local maxima --- Gribov copies~\cite{Gri,Zwa} ---
means P6-F1/F2 may select different representatives on different runs,
introducing an irreducible systematic. This section quantifies that
uncertainty for $\Gcross$.

\subsection{The Gribov Problem in $\SU(2)$ Lattice Landau Gauge}
\label{ssec:gribovproblem}

Stationarity of $\FL$ ($\Theta=0$) is necessary but not sufficient for
uniqueness. Gribov~\cite{Gri} showed multiple stationary points exist on
every orbit; Singer~\cite{Singer} proved a topological no-go theorem for
global gauge conditions over a compact base. The first Gribov region
$\Omega$ consists of Landau-gauge configurations with positive-definite
Faddeev--Popov operator; the fundamental modular region
$\Lambda\subset\Omega$ of global maxima underlies the
Gribov--Zwanziger programme~\cite{Zwa}, beyond our scope. In practice
one adopts either first-copy or best-copy prescriptions;
\cite{CM} quantifies the difference for the $\SU(2)$ gluon propagator
(and \cite{Bo} for $\SU(3)$).

Because $\Gcross$ is not gauge invariant (gauge non-invariance, Paper~5 \S8.1),
distinct copies give distinct values
$\Gcross^{(k)}(t)=8\sum_{x,\mu}\Re[b^{(k)}_t\bar b^{(k)}_0]$. This
copy-dependence must be characterised before $\Gcrossgf$ can serve as a
physical observable. SL5 was conditionally lifted by Landau-gauge
convergence (D19--D21); its full discharge additionally requires the
Gribov systematic to be bounded, which Proposition 7 provides.

\subsection{Best-Copy Strategy P6-F3}
\label{ssec:F3}

\paragraph{Algorithm P6-F3.}
For each configuration: (1)~generate $n_{\mathrm{start}}\ge5$
independent Haar-random gauge transformations $\{R^{(k)}(x)\}$;
(2)~apply P6-F1 ($N=6$) or P6-F2 ($N\ge8$) from each start, reaching
$\Theta^{(k)}<\epsg$; (3)~record $\FL^{(k)}$ and $\Gcross^{(k)}(t)$;
(4)~select the best copy $k^{*}=\operatorname{argmax}_k\FL^{(k)}$;
(5)~use $\{U^{\mathrm{best}}\}$ for all subsequent measurements. We
write $\Gcross^{\mathrm{best}}$, $\Gcross^{\mathrm{first}}$,
$\Delta\Gcross(t)=\max_k\Gcross^{(k)}(t)-\min_k\Gcross^{(k)}(t)$, and
$\Delta_{\mathrm{rel}}(t)=\Delta\Gcross(t)/|\Gcross^{\mathrm{best}}(t)|$.

\subsection{Proposition 7: Bounds on Gribov Copy Variation}
\label{ssec:prop7}

\setcounter{proposition}{6}
\begin{proposition}[Bounds on Gribov Copy Variation of $\Gcross$]
Let $\{U^{(k)}_\mu(x)\}$, $k=1,\ldots,n_{\mathrm{start}}\ge5$, be Gribov
copies of the same configuration produced by P6-F3, all satisfying the
Landau condition to $\epsg=10^{-14}$.
\begin{enumerate}[label=(\roman*)]
\item \textbf{Algebraic variation bound.} For any two copies $k,l$,
\begin{equation}
  |\Gcross^{(k)}(t)-\Gcross^{(l)}(t)|\le 8dV(\delta_t+\delta_0),
  \label{eq:P7bound}
\end{equation}
with $\delta_t=\max_{x,\mu}|b^{(k)}_t-b^{(l)}_t|$,
$\delta_0=\max_{x,\mu}|b^{(k)}_0-b^{(l)}_0|$. The bound is not saturated
on thermalised ensembles; the relative variation satisfies
$\Delta_{\mathrm{rel}}=O(1)$ in the coupling.
\item \textbf{Best-copy improvement.} Since $\Gcross$ is not gauge
invariant, $\Gcross^{\mathrm{best}}\ne\Gcross^{\mathrm{first}}$
generically; best-copy selection reduces fluctuations (driving links
toward $\Id_2$, $b=0$) but imposes no definite ordering on $\Gcross$.
\item \textbf{$\be$-dependent threshold (empirical, D22).} Three regimes:
(a)~$\be\ge9.0$, $N=6$: $\Delta_{\mathrm{rel}}<1\%$ (manageable);
(b)~$\be\in\{5,6\}$, $N=6$: $\Delta_{\mathrm{rel}}\in[1\%,10\%]$
(moderate); (c)~$\be\le3.5$, extended ensembles:
$\Delta_{\mathrm{rel}}>10\%$ (severe).
\end{enumerate}
\end{proposition}

\begin{proof}
\textit{(i)} By the triangle inequality and adding/subtracting
$b^{(l)}_t\bar b^{(k)}_0$,
$|b^{(k)}_t\bar b^{(k)}_0-b^{(l)}_t\bar b^{(l)}_0|\le
|b^{(k)}_t-b^{(l)}_t||\bar b^{(k)}_0|+|b^{(l)}_t||\bar b^{(k)}_0-\bar
b^{(l)}_0|$. Since $|b|\le1$ by unitarity, summing over $dV$ links gives
\eqref{eq:P7bound}; equality would require $|b|=1$ everywhere, excluded
by Landau gauge. By the $b$-transformation lemma (Paper 5, \S8.1), at weak coupling copies differ by
near-identity transformations, so $\delta_t=O(g^2)$ and
$|\Gcross^{(k)}-\Gcross^{(l)}|=O(g^2\cdot4V)$. Then
$\Delta_{\mathrm{rel}}=O(g^2)/(g^2/2)=O(1)$: the $O(g^2)$ factors cancel,
leaving the variation parametrically unsuppressed and controllable only
empirically.

\textit{(ii)} By the gauge non-invariance result (Paper~5 \S8.1), distinct gauge-fixed
configurations of one orbit give distinct values; $g^{\mathrm{best}}$ and
$g^{\mathrm{first}}$ differ (D22), so the values differ generically.
Maximising $\FL$ minimises the average $|b|^2$ per link, reducing
sensitivity, but fixes no sign of the difference.

\textit{(iii)} Empirical; confirmed by D22. The percentage thresholds are
order-of-magnitude expectations replaced by measured values in the
verification summary.
\end{proof}

\paragraph{SL5 discharge.}
With Proposition 7 established and D22--D23 passed, SL5 is fully
discharged: the non-perturbative identification
$\Gcrossgf\leftrightarrow W^{\pm}$ (Theorem 2) rests on verified Landau
convergence (D19--D21) and a characterised Gribov systematic.

\subsection{Statistics and Reporting Conventions}
\label{ssec:gribovstats}

For each ensemble and timeslice we report
$\langle\Gcross^{\mathrm{best}}\rangle(t)$, the statistical error
$\sigma_{\mathrm{stat}}(t)$, the mean copy range
$\langle\Delta\Gcross\rangle(t)$, its spread
$\sigma_{\mathrm{Gribov}}(t)$, and $\Delta_{\mathrm{rel}}(t)$. The ratio
$\sigma_{\mathrm{Gribov}}/\sigma_{\mathrm{stat}}$ dictates reporting: if
$<0.1$, statistical error only (Gribov controlled); if $\in[0.1,1]$,
both reported; if $>1$, Gribov dominant and the ensemble is flagged. The
first-copy value $\Gcross^{\mathrm{first}}$ is recorded alongside for
like-for-like comparison with the first-copy $\SU(2)$ data of \cite{CM}.

\paragraph{Checks D22 and D23.}
D22 (inverted) requires measurably different copies:
$\max_{t,\mathrm{cfg}}\Delta\Gcross(t)>\varepsilon_{\mathrm{copy}}=10^{-6}$.
The inversion is motivated by the expectation that thermalised
non-Abelian configurations generically admit multiple Gribov
copies~\cite{Gri}; Singer~\cite{Singer} establishes the continuum
obstruction to a global gauge condition but does not by itself guarantee
copy multiplicity for a given finite-lattice configuration. Persistent
degeneracy across independent random starts is therefore flagged for a
direct check of start independence (RNG and initialisation), not assumed
to be a code defect; in the perturbative limit $\be=20.0$ near-degeneracy
is physical (``copy-degenerate regime''). D23
verifies best-copy stability under a small random perturbation and
re-fixing, with failure rates above $5\%$ signalling Gribov-horizon
proximity.

%=======================================================================
\section{The Gauge-Fixed \texorpdfstring{$W^{\pm}$}{W} Correlator:
Theorem 2}
\label{sec:theorem2}

With the ensemble in a controlled gauge (D19--D20) and Gribov ambiguity
handled by P6-F3 (D23), we define the gauge-fixed correlator, derive its
weak-coupling relation to the $W^{\pm}$ field, and establish Theorem 2
with a three-term error budget. The scope of this section is the Paper
4/5 ensembles ($N=6$, $\be\in\{5,6,9,12\}$); extended ensembles are
introduced in \Cref{sec:ensembles}.

\subsection{The Gauge-Fixed Observable (P6-F4)}
\label{ssec:gfobs}

\begin{definition}[Gauge-Fixed Cross-Channel Correlator]
\begin{equation}
  \Gcrossgf(t)=8\sum_{x,\mu}\Re\bigl[b^{\mathrm{gf}}_t(x,\mu)\,
  \bar b^{\mathrm{gf}}_0(x,\mu)\bigr],
  \label{eq:gcrossgf}
\end{equation}
where $b^{\mathrm{gf}}(x,\mu)=(U^{g^{\mathrm{best}}}_\mu(x))_{01}$ is the
$(0,1)$ element of the best-copy gauge-fixed link from P6-F3.
\end{definition}

The ensemble average $\langle\Gcrossgf(t)\rangle$ is the mean over
configurations. Reproducibility under repeated measurement (deterministic
arithmetic) is verified by D24. The gauge-fixing correction
$\Delta^{\mathrm{gf}}(t;c)=\Gcrossgf(t;c)-\Gcross(t;c)$ is generically
non-zero since $\Gcross$ is gauge dependent.

\subsection{Weak-Coupling Identification}
\label{ssec:weakcoupling}

For $A=A^a\tau^a$ traceless antihermitian, the $\SU(2)$ Cayley formula
gives $\exp(igA)=\cos(g|A|/2)\Id+(2i/|A|)\sin(g|A|/2)A$, so
\begin{equation}
  b=(\exp(igA))_{01}=(ig/\sqrt2)\,W
  \bigl[1-g^2|A|^2/24+O(g^4)\bigr],
  \label{eq:bexpand}
\end{equation}
with $W\equiv(A^1-iA^2)/\sqrt2$. Because $A^2=(|A|^2/4)\Id$ for
traceless antihermitian $2\times2$ $A$, there is no $O(g^2)$ correction
to the off-diagonal element:
\begin{equation}
  b=(ig/\sqrt2)\,W+O(g^3).
  \label{eq:bclean}
\end{equation}
This $\SU(2)$ improvement tightens the Theorem 2 error from $O(g)$ to
$O(g^2)$ relative. Conjugating ($ig\to-ig$),
\begin{equation}
  \Re[b_t\bar b_0]=\tfrac{g^2}{4}(A^1_tA^1_0+A^2_tA^2_0)+O(g^4),
  \label{eq:Rebb}
\end{equation}
the positive sign arising from $(ig)(-ig)=+g^2$ (the locked Paper 5
Q5.3 sign). Hence
\begin{equation}
  \Gcrossgf(t)=2g^2\sum_{x,\mu}(A^1_tA^1_0+A^2_tA^2_0)+O(g^4).
  \label{eq:gcrossOg2}
\end{equation}

\subsection{Theorem 2}
\label{ssec:theorem2statement}

\setcounter{theorem}{1}
\begin{theorem}[Gauge-Fixed $W^{\pm}$ Propagator; paper-local]
In Landau gauge (Lemma 4), on best-copy gauge-fixed configurations
(Proposition 7, P6-F3), the zero-spatial-momentum $W^{\pm}$ temporal
propagator satisfies
\begin{equation}
  \DW(t)=\frac{\Gcrossgf(t)}{8Vg^2}
  +\delta_{\mathrm{pert}}+\delta_{\mathrm{Gribov}}+\delta_{\mathrm{stat}},
  \label{eq:thm2}
\end{equation}
where
\begin{equation}
  \DW(t)\equiv\frac{1}{4V}\sum_{x,\mu}
  \bigl(\langle A^1_\mu(x,t)A^1_\mu(x,0)\rangle
  +\langle A^2_\mu(x,t)A^2_\mu(x,0)\rangle\bigr).
  \label{eq:DWdef}
\end{equation}
The leading-order identification $\DW(t)=\Gcrossgf(t)/(8Vg^2)+O(g^2)D^{(0)}$
holds with $O(g^2)$ relative error (not $O(g)$), by the $\SU(2)$ group
identity of \S\ref{ssec:weakcoupling}. Using $g^2=(16/3)(1-\langle
P\rangle)=8/\be+O(\be^{-2})$,
\begin{equation}
  \DW(t)=\Gcrossgf(t)\,\frac{\be}{64V}+O(\be^{-1}).
  \label{eq:working}
\end{equation}
L/R caveat: the embedding does not, by itself, distinguish $\SU(2)_L$
from $\SU(2)_R$; this requires external fermionic input.
\end{theorem}

\begin{proof}
\textit{Step 1 (gauge-fixing validity).} By Lemma 4, P6-F1/F2 converge to
$\epsg=10^{-14}$ (D19); the per-site condition holds (D20); by
Proposition 7 and P6-F3 the best copy is selected and stable (D23), so
$\Gcrossgf$ is well defined and reproducible.
\textit{Step 2 (expansion).} From \eqref{eq:bclean}--\eqref{eq:Rebb} the
relative perturbative error is $O(g^2)$.
\textit{Step 3 (identification).} From \eqref{eq:gcrossOg2} and
\eqref{eq:DWdef}, $8Vg^2\DW(t)=2g^2\sum_{x,\mu}(A^1A^1+A^2A^2)$ exactly,
so $\Gcrossgf(t)-8Vg^2\DW(t)=O(g^4)$ and the stated identification
follows.
\textit{Step 4 (coupling).} Using $g^2=(16/3)(1-\langle P\rangle)$ introduces
$O(g^2)$ relative error, consistent with Step 2.
\textit{Step 5 (remainders).} $\delta_{\mathrm{Gribov}}$ is bounded by
Proposition 7 (D22--D23); $\delta_{\mathrm{stat}}\sim1/\sqrt{\Ncfg}$.
\end{proof}

\begin{remark}[Provenance]
The Paper 5 \S8.3 identification was perturbative only, not gauge-fixed,
without Gribov analysis, and unverified. Theorem 2 supplies all four,
plus the $O(g^2)$ (rather than $O(g)$) error bound from the $\SU(2)$
identity. It is a new result of Paper 6.
\end{remark}

\subsection{Error Budget}
\label{ssec:errorbudget}

With $D^{(0)}(t)\equiv\Gcrossgf(t)/(8Vg^2)$, the three independent terms
are: the perturbative correction
$|\delta_{\mathrm{pert}}|/D^{(0)}\le C_P g^2\approx8C_P/\be$ (so $\sim40\%
C_P$ at $\be=20$, $O(1)$ for $\be\lesssim8$); the Gribov systematic
$\delta_{\mathrm{Gribov}}(t)=\Delta\Gcrossgf{}^{\max}(t)/(8Vg^2)$ from
D22--D23; and the statistical uncertainty
$\sigma_{\mathrm{stat}}[\DW(t)]=\sigma[\Gcrossgf(t)]/(8Vg^2\sqrt{\Ncfg})$.
Combined in quadrature,
\begin{equation}
  [\delta\DW(t)]^2=\delta_{\mathrm{pert}}^2+\delta_{\mathrm{Gribov}}^2
  +\delta_{\mathrm{stat}}^2.
  \label{eq:budget}
\end{equation}

\subsection{Corollary 6: Sign Structure of \texorpdfstring{$\Gcrossgf$}{Gcross}}
\label{ssec:cor6}

\setcounter{corollary}{5}
\begin{corollary}[Sign Structure of $\Gcrossgf$; series-wide]
\hfill
\begin{enumerate}[label=(\roman*)]
\item At leading order $O(g^2)$, the ensemble average satisfies
$\langle\Gcrossgf(t)\rangle\ge0$ for all $t\ge0$. This is an
ensemble-average statement, not pointwise: individual configurations may
be negative at finite volume.
\item At non-perturbative level, $\langle\Gcrossgf(t)\rangle$ is not
sign-definite once $O(g^4)$ and genuine non-perturbative contributions
dominate. A threshold $\bc$ is defined as the coupling below which the
average turns negative for some $t$.
\item The sign transition at $\bc$ is a diagnostic for the onset of
non-perturbative dynamics in the $W^{\pm}$ channel; its $\be$- and
$N$-dependence are the empirical subject of \Cref{sec:signstructure}.
\end{enumerate}
\end{corollary}

\begin{proof}
\textit{(i)} From \eqref{eq:gcrossOg2},
$\langle\Gcrossgf(t)\rangle=2g^2\sum_{x,\mu}\langle
A^1_tA^1_0+A^2_tA^2_0\rangle+O(g^4)$; at strict $O(g^2)$ the right-hand
side is the \emph{free} lattice gauge-field two-point function, which is
reflection positive, so $\langle\Gcrossgf(t)\rangle\ge0$ for $t\ge0$. (This
is a free-theory statement: the interacting Landau-gauge gluon propagator
is \emph{not} reflection positive, which is precisely the non-perturbative
sign violation of part~(ii) and \Cref{sec:signstructure}.)
\textit{(ii)} By example at $\be=2.0$, non-perturbative corrections are
$O(1)$ and negative-valued configurations are present.
\end{proof}

\begin{remark}[Correction note]
The original task specification stated $\Gcrossgf\ge0$ pointwise; this is
too strong at finite volume. The correct statement is the ensemble
average of part~(i).
\end{remark}

\subsection{Checks D24 and D25}
\label{ssec:D24D25}

D24 (reproducibility) verifies
$\max_{t,c}|G_1(t;c)-G_2(t;c)|<\varepsilon_{D24}=10^{-12}$; a
floating-point-floor result ($\sim10^{-14}$) confirms purely
deterministic arithmetic. D25 (weak-coupling consistency) at
$\be\in\{12,20\}$ requires (a)~ensemble-average positivity, (b)~$\cosh$-like
(periodic) behaviour of $\DW(t)$, symmetric about $t=\Nt/2$, within
$2\sigma$, (c)~consistency at $\be=20$ with the near-massless free-mode
form: the zero-momentum correlator is there dominated by the free mode, so
a stable $m_{\mathrm{eff}}>0$ plateau is a non-perturbative (not a
weak-coupling) expectation and is not required here, and (d)~coupling
consistency $|D^{\mathrm{meas}}_{W}-\DW|<
(16/\be^2)\DW$. The $\be=20$ component is conditional on the extended
ensemble (\Cref{sec:ensembles}).

%=======================================================================
\section{Ensemble Extension and Volume Scaling}
\label{sec:ensembles}

The base ensembles ($N=6$, $\Ncfg=20$) carried two limitations: a $\sim22\%$
statistical uncertainty and an unremovable finite-volume systematic from a
single lattice size. Both are eliminated here.

\subsection{Rationale and Specifications}
\label{ssec:ensemblespec}

Raising $\Ncfg$ from $20$ to $100$ reduces the fractional statistical
error by $\sqrt5\approx2.2$, below $10\%$ for $\be\ge5$. Extending to
$N\in\{8,12,16\}$ exposes the $1/N^2$ finite-volume scaling and enables
the thermodynamic extrapolation required for comparison with the
$\SU(2)$ data of \cite{CM}.
The base ensembles ($N=6$, $\be\in\{5,6,9,12\}$) are reprocessed with
$\Ncfg=100$. Three groups of new ensembles are generated (Wilson plaquette
action, heatbath plus overrelaxation, periodic boundary conditions,
$\ge500$ thermalisation sweeps, $\ge10$ sweeps between measurements):

\begin{center}
\renewcommand{\arraystretch}{1.15}
\begin{tabular}{lll}
\toprule
Group & Coverage & Count\\
\midrule
E1 (strong coupling) & $N\in\{6,8,12,16\}$, $\be\in\{2.0,2.5,3.0,3.5\}$ & 16\\
E2 (perturbative)    & $N\in\{6,8,12,16\}$, $\be=20.0$ & 4\\
E3 (FV scaling)      & $N\in\{8,12,16\}$, $\be\in\{5,6,9,12\}$ & 12\\
\bottomrule
\end{tabular}
\end{center}

\noindent
giving 32 new ensembles, all $\Ncfg=100$. Group E2 discharges FLAG-7 from
\Cref{sec:theorem2} (the $N=6$, $\be=20$ ensemble required by D25).
Algorithm selection follows Lemma 4: P6-F1 at $\be\le9$ or $N\le8$,
P6-F2 at $\be\ge9$ or $N\ge12$.

\subsection{Autocorrelation and Jackknife Errors}
\label{ssec:jackknife}

The integrated autocorrelation time
$\tau_{\mathrm{int}}=\tfrac12+\sum_{\tau\ge1}\rho(\tau)$ is estimated on
$\Gcrossgf(t=1)$ with the Madras--Sokal window. Configurations are
independent when $n_{\mathrm{sep}}\ge2\tau_{\mathrm{int}}$; the
provisional $n_{\mathrm{sep}}=10$ is increased where
$\tau_{\mathrm{int}}>5$. All errors in \Cref{sec:ensembles,sec:comparison,sec:signstructure}
are jackknife errors,
\begin{equation}
  \sigma^2_{\mathrm{JK}}(t)=\frac{\Ncfg-1}{\Ncfg}
  \sum_c\bigl(\bar G_{-c}(t)-\bar G(t)\bigr)^2,
  \label{eq:jackknife}
\end{equation}
computed on configuration-level data binned into blocks of length
$\ge2\tau_{\mathrm{int}}$ (block jackknife), so that residual
autocorrelation is captured, $\sigma^2_{\mathrm{JK}}/\sigma^2_{\mathrm{naive}}
\approx2\tau_{\mathrm{int}}$ for the binned estimator. A delete-one
jackknife on unbinned, autocorrelated configurations would instead return
$\approx\sigma^2_{\mathrm{naive}}$ and \emph{underestimate} the error; the
bin length is increased until $\sigma_{\mathrm{JK}}$ plateaus.

\paragraph{Check D26.}
With coefficient of variation
$\mathrm{CV}(t{=}1)=\sigma_{\mathrm{JK}}(t{=}1)/|\langle\Gcrossgf(t{=}1)\rangle|$,
the pass condition is $\mathrm{CV}<5\%$ for all ensembles plus
running-mean stability. PARTIAL PASS is expected at strong coupling
($\be\le3.5$), where non-perturbative fluctuations inflate the variance.

\subsection{Finite-Volume Scaling}
\label{ssec:FV}

The zero-momentum propagator on an $N^3$ spatial lattice is fitted with the
phenomenological finite-volume ansatz
\begin{equation}
  \DW(t;N)=\DW^{\infty}(t)+\frac{a(t)}{N^2}+O(N^{-4}),
  \label{eq:FVansatz}
\end{equation}
with $a(t)$ a fit parameter. The $1/N^2$ form is \emph{not} obtained from
Euler--Maclaurin: for a smooth periodic summand on a full-period grid the
Euler--Maclaurin boundary terms cancel and the residual error is
super-algebraically small, so any power-law correction must originate
elsewhere --- in the long-distance (near-massless / infrared)
non-analyticity of the two-point function and in the treatment of the
spatial zero mode. The $1/N^2$ leading power is therefore adopted
empirically and validated by the fit. For a decoupling (massive,
$m_{\mathrm{IR}}>0$) propagator the leading finite-volume correction is
instead expected to be exponential, $\sim e^{-m_{\mathrm{IR}}N}$; both the
$1/N^2$ and an $e^{-mN}$ form are fitted and the better-supported one
retained, with the $N^{-4}$ term monitored at small $\be$. The ansatz
presumes $m_{\mathrm{IR}}N\gg1$ (expected for $\be\ge5$). For each
$(t,\be)$, \eqref{eq:FVansatz} is fitted across $N\in\{6,8,12,16\}$ by
weighted least squares ($\chi^2/\mathrm{dof}=\chi^2/2$, target $\lesssim1$
at $\be\ge5$), yielding $\DW^{\infty}(t)$ with propagated jackknife
uncertainty and a residual systematic $\delta_{\mathrm{FV}}(t)=|a(t)|/256$.

\subsection{Plaquette Coupling}
\label{ssec:coupling}

The working formula \eqref{eq:working} requires $g^2$ per ensemble,
extracted as $g^2_{\mathrm{meas}}=(16/3)(1-\langle P\rangle)$ (Table B.1). At
$\be=20$ the deviation from $g^2_{\mathrm{pert}}=8/\be$ is a higher-order
effect to be quantified from the measured $\langle P\rangle$; a deviation
exceeding the one-loop estimate triggers substitution of
the one-loop improved coupling. At strong coupling, large deviations are
expected and reflect genuine non-perturbative physics, not failure.

%=======================================================================
\section{Validation Against the $\SU(2)$ Landau-Gauge Gluon Propagator}
\label{sec:comparison}

We ask whether $\DW(t)$, extracted from the $b$-channel, agrees with
independently measured $\SU(2)$ Landau-gauge gluon propagator data. This
requires colour symmetry (\S\ref{ssec:colour}), normalisation matching
(\S\ref{ssec:norm}), and the absence of dominant uncontrolled systematics.

\subsection{The Comparison Target}
\label{ssec:target}

Cucchieri--Mendes~\cite{CM} measured the zero-momentum $\SU(2)$
Landau-gauge gluon propagator on large lattices (to $N=128$,
$\beta\in\{2.2,2.4\}$), averaging over all three $\SU(2)$ colour components
and all directions, and found $D(p^2{=}0)$ finite and non-zero with strong
infrared suppression --- the decoupling solution,
$D(p^2)\sim 1/(p^2+m_{\mathrm{IR}}^2)$. This is our $\SU(2)$ benchmark.
Bogolubsky et al.~\cite{Bo} report the same decoupling behaviour but in
$\SU(3)$ gluodynamics (lattices to $L=96$ at $\beta=5.70$, with a
simulated-annealing approach to the global gauge-functional maximum); it is
cited here only as a qualitative, gauge-group-independent confirmation of
the decoupling scenario and a precedent for global-maximum gauge fixing,
\emph{not} as an $\SU(2)$ quantitative benchmark.
\begin{quote}\itshape\textbf{[OPEN]} The matched-$\be$ Part~B comparison
(\S\ref{ssec:scopesplit}) requires a genuine large-volume $\SU(2)$ dataset.
A suitable reference is Bornyakov, Mitrjushkin and M\"uller-Preussker,
\emph{Phys.\ Rev.\ D} \textbf{81} (2010) 054503 ($\SU(2)$ Landau-gauge
gluon propagator; continuum and infinite-volume limits, infrared mass
scale). It must be added to the bibliography and used in place of \cite{Bo}
for any quantitative $\SU(2)$ comparison.\end{quote}

\subsection{$\SU(2)$ Colour Symmetry}
\label{ssec:colour}

\begin{proposition}[Colour Symmetry of the Landau-Gauge Propagator]
In $\SU(2)$ lattice gauge theory in Landau gauge,
\begin{equation}
  D^{a=1}(p^2)=D^{a=2}(p^2)=D^{a=3}(p^2)\equiv D(p^2).
  \label{eq:colour}
\end{equation}
\end{proposition}

\begin{proof}
The Wilson action is invariant under the global adjoint action
$A^a_\mu\to R^{ab}A^b_\mu$, $R\in\SO(3)$. Since $\SO(3)$ acts
transitively on $\{1,2,3\}$, for any $a,b$ there is $R$ with
$R\hat e_a=\hat e_b$; the Landau condition and $\FL$ are invariant, so
$\langle(A^a_\mu)^2\rangle=\langle(A^b_\mu)^2\rangle$, giving
\eqref{eq:colour}. (Transitivity of $\SO(3)$ on the adjoint index is
specific to $\SU(2)$.)
\end{proof}

\noindent
Consequently $\DW(p^2)=D^{a=1}(p^2)=D(p^2)$: the $b$-channel propagator
equals the three-colour average used by \cite{CM}, so the comparison is
like-for-like. The $W^3/A^3$ component (measured by $\Gdiag$) must equal
$\DW$ by \eqref{eq:colour}, providing an independent cross-check
(S6-2, \S\ref{ssec:S62}); no $\SU(2)_L/\SU(2)_R$ distinction is involved.

\paragraph{Status of this comparison.}
Equation \eqref{eq:colour} fixes how this section should be read: in the
$\SU(2)$ theory the $b$-channel propagator $\DW$ \emph{is} the standard
Landau-gauge gluon propagator, not a distinct observable. The novelty of
the $b$-channel construction is therefore \emph{structural} --- $\DW$ is
definable, and cleanly isolated, within the $\Cl(0,2)$ framework --- rather
than \emph{informational}: it carries no physical content that the standard
gluon propagator does not already express. Agreement with the reference
data is consequently \emph{expected} by colour symmetry, so this section is
a validation and implementation cross-check of the $\Cl(0,2)$ extraction,
not a new physical determination. A genuinely new propagator would require
breaking the colour symmetry of \eqref{eq:colour} --- e.g.\ through the
fermionic $\SU(2)_L/\SU(2)_R$ content of \Cref{sec:fermionic} --- which is
beyond the pure-gauge scope of this section.

\subsection{Normalisation Matching}
\label{ssec:norm}

With $\DW(t)$ as in \eqref{eq:DWdef}, the temporal DFT gives
$\DWt(\omega{=}0)=\sum_{t}\DW(t)$. Converting to the \cite{CM}
convention requires tracking four factors. (i)~\emph{Colour.} By
\eqref{eq:colour}, $\DW=(1/4V)\sum_{x,\mu}(\langle A^1A^1\rangle+\langle
A^2A^2\rangle)=(1/2V)\sum_{x,\mu}\langle A^1A^1\rangle$ is a
two-charged-colour \emph{sum}, whereas \cite{CM} uses the three-colour
\emph{average} $D=\tfrac13\sum_a D^a$; with $D^1=D^2=D^3$ this is a
relative factor $\tfrac12$ (our sum is twice their average).
(ii)~\emph{Direction.} A factor $1/d$ converts our direction sum to their
per-direction average. (iii)~\emph{Field convention.} The complex-field
relation $C_{W^\pm}=2\DW$ (Paper 5) contributes a further factor that must
be confirmed against the Paper 5 definition. (iv)~\emph{DFT/volume.} The
$1/\Nt$ DFT normalisation, with $V$ as fixed by the Papers 4--5 convention
(spatial $N^3$ vs total $N^4$). Collecting these,
\begin{equation}
  D_{\mathrm{CM}}(p^2{=}0)
  =\kappa_{\mathrm{CM}}\,
   \frac{\sum_t\Gcrossgf(t)\,\be/(64V)}{d\,\Nt},
  \label{eq:DCM}
\end{equation}
where $\kappa_{\mathrm{CM}}=\tfrac12\,c_{\mathrm{conv}}$ absorbs the colour
factor $\tfrac12$ of (i) together with the field-convention and zero-mode
constants of (iii)--(iv).
\begin{quote}\itshape\textbf{[OPEN]} $\kappa_{\mathrm{CM}}$ is not pinned in
this draft. It must be derived from the verbatim \cite{CM} zero-mode
definition $\hat A^a_\mu(p{=}0)$ and the inherited Papers 4--5 $V$
convention before any quantitative comparison; the colour factor $\tfrac12$
is established here, the remainder is deferred.\end{quote}
The sanity check S6-3 verifies the \emph{implemented} conversion by
comparing the DFT of $\DW(t)$ to the direct zero-mode computation
\eqref{eq:DCM}, requiring their ratio to equal the implemented
$d\,\Nt/\kappa_{\mathrm{CM}}$ within $10^{-10}$.

\subsection{Scope Limitation and Two-Part Structure}
\label{ssec:scopesplit}

Since \cite{CM} works at $\beta\in\{2.2,2.4\}$ (our Scope D,
ungenerated), and our available ensembles have $\be\ge5$, the comparison
splits. \textbf{Part A (available):} at $\be\in\{5,6,9,12,20\}$,
$D_{\mathrm{CM}}(p^2{=}0)$ is computed via \eqref{eq:DCM} and its
$\be$-dependence examined as an internal consistency check of Theorem 2.
\textbf{Part B (deferred):} at $\be\in\{2.0,2.5,3.0,3.5\}$, the
FV-extrapolated $\DW^{\infty}(\omega{=}0)$ at $N=16$ is compared directly
to the $\SU(2)$ data of \cite{CM} (see the \S\ref{ssec:target} [OPEN] note
on adding a dedicated large-volume $\SU(2)$ reference); this requires Scope D
generation and is fully specified
here for later execution.

\subsection{$\be$-Dependence and the Onset of Suppression}
\label{ssec:partA}

At weak coupling, $\Gcrossgf\propto g^2\propto\be^{-1}$, so that
$\DW=\Gcrossgf\,\be/(64V)$ is $\be$-independent at leading order; summing
timeslices (cf.\ \S\ref{ssec:plateau}),
\begin{equation}
  \DWt(\omega{=}0)=(d/2)\,\widetilde C_{\mathrm{free}}(\omega{=}0)+O(\be^{-1})
  \equiv P_\infty+O(\be^{-1}),
  \label{eq:perturbgrowth}
\end{equation}
a $\be$-independent plateau at leading order, with a sub-leading slope that
is power-law (not exponential) in $\be^{-1}$. The \cite{CM,Bo} infrared
suppression is a non-perturbative effect visible only at $\be\sim2$--$4$,
where $\DWt(\omega{=}0)$ deviates \emph{downward} from the plateau $P_\infty$.
Part A ensembles ($\be\ge5$) sit on the plateau; the onset of this downward
deviation as $\be$ decreases is the signal sought. At $\be=20$,
$D_{\mathrm{CM}}(p^2{=}0)$ should match the plateau \eqref{eq:perturbgrowth};
deviations beyond the \Cref{sec:theorem2} budget would indicate an
implementation error.

\subsection{Systematic Budget and Cross-Checks}
\label{ssec:S62}

The total systematic on $\DWt(\omega{=}0)$ combines perturbative, Gribov,
finite-volume, and statistical terms in quadrature, paralleling
\eqref{eq:budget}. The dominant term is $\delta_{\mathrm{stat}}$ or
$\delta_{\mathrm{FV}}$ at $\be\ge9$, $\delta_{\mathrm{pert}}$ near
$\be=5$, and $\delta_{\mathrm{Gribov}}$ at $\be\le3.5$ (limiting Part B).
The colour-symmetry cross-check S6-2 requires the normalised difference
$|D_{\mathrm{diag}}(\omega{=}0)-\DW(\omega{=}0)|/\sigma_{\Delta}<3$, where
$\sigma_{\Delta}$ is the jackknife error of the per-configuration difference
$D_{\mathrm{diag}}-\DW$ (not of either term separately, since the two
channels are strongly correlated on shared configurations),
for all available ensembles; agreement validates simultaneously the
$b$-channel extraction, the $\Cl(0,2)$ decomposition, and the absence of
colour-asymmetric gauge-fixing bias. Future requirement F9 is specified:
Part A is runnable on available ensembles (pending execution); Part B is
specified pending Scope D.

%=======================================================================
\section{Sign Structure and \texorpdfstring{$\be$}{beta} Characterisation}
\label{sec:signstructure}

This section delivers Corollary 6 Part~(iii): the empirical
characterisation of the threshold $\bc$ below which
$\langle\Gcrossgf(t)\rangle$ loses sign-definiteness, closing Open
Question OQ2 of Paper 5. No new numbered results are introduced; the
content is computational. The observable is $\DW(t)$ (Theorem 2);
Observable 2 does not appear.

\begin{quote}\itshape
L/R caveat: The embedding does not, by itself, distinguish $\SU(2)_L$
from $\SU(2)_R$; this requires external fermionic input (\Cref{sec:fermionic}).
\end{quote}

\subsection{Sign Analysis Protocol (P6-F6, P6-F7)}
\label{ssec:signprotocol}

For each ensemble the jackknife estimator yields
$\langle\Gcrossgf(t)\rangle\pm\sigma_{\mathrm{JK}}(t)$ for
$t=1,\ldots,\Nt/2$. Timeslices are classified, with a uniform $2\sigma$
threshold, as
\begin{align}
  \text{P (Positive):}&\quad \langle\Gcrossgf(t)\rangle>+2\sigma_{\mathrm{JK}}(t),\notag\\
  \text{N (Negative):}&\quad \langle\Gcrossgf(t)\rangle<-2\sigma_{\mathrm{JK}}(t),\\
  \text{Z (Zero):}&\quad |\langle\Gcrossgf(t)\rangle|\le 2\sigma_{\mathrm{JK}}(t).\notag
\end{align}
Ensembles are typed Type-P (all P), Type-N (some N, none P), Type-M
(both P and N), or Type-Z (all Z). Type-P is expected at weak coupling
(Corollary 6(i)); Type-N/M is the non-perturbative signal; Type-Z marks
the crossover. The classification is applied to both $\Gcrossgf$ and
$\Gdiaggf$ (the latter defined via the $a$-element by P6-F6).

\subsection{Threshold $\bc$ Characterisation}
\label{ssec:betac}

\begin{equation}
  \bc(N)=\sup\{\be:\text{ensemble type}\in\{N,M,Z\}\}.
  \label{eq:betacdef}
\end{equation}
If all available ensembles at $N$ are Type-P, then $\bc(N)<5.0$ pending
Scope D --- a strict upper bound, not a failure. When the transition is
bracketed by adjacent couplings, $\bc(N)$ is estimated by linear
interpolation in the Category-P fraction. If $\bc(N)$ is obtained as a
float for $N\in\{6,8,12,16\}$, the $1/N^2$ ansatz of
\Cref{sec:ensembles} is applied,
\begin{equation}
  \bc(N)=\bc^{\infty}+A/N^2+O(N^{-4}),\qquad\chi^2/\mathrm{dof}<3.
  \label{eq:betacfit}
\end{equation}
If $\chi^2/\mathrm{dof}\ge3$ (FLAG-8), $\bc^{\infty}$ is reported as a
lower bound. If $\bc^{\infty}$ is consistent with zero within
$2\sigma_{\mathrm{fit}}$, this is flagged as a significant secondary
finding (sign-positivity in the thermodynamic limit), deferred pending
Scope D.

\subsection{Colour Symmetry Cross-Check}
\label{ssec:colourcrosscheck}

The diagonal correlator
$\Gdiaggf(t)=8\sum_{x,\mu}\Re[a^{\mathrm{gf}}_t\bar a^{\mathrm{gf}}_0]$
is classified by the identical protocol. By colour symmetry
(\S\ref{ssec:colour}), the sign pattern of $\Gdiaggf$ should agree
ensemble-by-ensemble with $\Gcrossgf$ at weak coupling; discrepancies are
reported without suppression as inline flags (candidate causes: lattice
artefacts, gauge-fixing residuals, finite-$N$ boundary effects).

\subsection{$\DWt(\omega{=}0)$ Plateau and Onset Coupling (P6-F8)}
\label{ssec:plateau}

With $\DWt(\omega{=}0)=\sum_t\DW(t)
=\sum_t\langle\Gcrossgf(t)\rangle\be/(64V)+O(\be^{-1})$, the
weak-coupling expansion gives $\DW(t)=(d/2)\bar C_{\mathrm{free}}(t)
+O(\be^{-1})$ with $\bar C_{\mathrm{free}}$ a $\be$-independent
lattice-geometry constant. Hence
\begin{equation}
  \DWt(\omega{=}0)=(d/2)\,\widetilde C_{\mathrm{free}}(\omega{=}0)
  +O(\be^{-1}),
  \label{eq:plateauval}
\end{equation}
flat in $\be$ at leading order --- a consequence of Theorem 2, with the
plateau value $P_\infty=(d/2)\widetilde C_{\mathrm{free}}(\omega{=}0)$.
The sub-leading slope is power-law in $\be^{-1}$ (not exponential). The
onset coupling is $\beta_{\mathrm{onset}}=\min\{\be:|\DWt(\omega{=}0;\be)
-P_\infty|<2\sigma_{\widetilde D}\}$, expected to agree with $\bc(N)$
within one $\be$ step. A downward deviation below $P_\infty$ at small
$\be$ is the decoupling signature of \cite{CM,Bo}.

\subsection{Check D27}
\label{ssec:D27}

D27 is a composite. \textbf{D27a (standard, runnable):} for
$N\in\{6,8,12,16\}$ and $\be\in\{9,12,20\}$,
$\langle\Gcrossgf(t)\rangle>0$ for all $t$ within $2\sigma_{\mathrm{JK}}$
(Type-P), confirming Corollary 6(i); twelve sub-conditions, PARTIAL PASS
on $N\in\{6,8\}$ if Scope C is ungenerated; any Category-N timeslice is a
FAIL requiring revisiting D19--D23. \textbf{D27b (inverted, deferred):}
at $\be=2.0$ (Scope D), $\langle\Gcrossgf(t)\rangle<-2\sigma_{\mathrm{JK}}$
for at least one timeslice per $N$ (PASS = negative sign observed). Per
QS5, a fully positive result at $\be=2.0$ is not a code error but implies
$\bc(N)<2.0$, reported as PARTIAL PASS with the bound. Composite: D27a
PASS with D27b DEFERRED gives D27 = PARTIAL PASS, the current expected
status.

%=======================================================================
\section{Fermionic Sector: Setup and the Cross-Projected Propagator}
\label{sec:fermionic}

This section establishes the fermionic groundwork by implementing the
Future Requirements F1--F4 of Paper 5 \S7.3 (paper-local functions
P6-F10--P6-F13). It does not perform the full non-perturbative fermionic
test, which is Paper 7. It (a)~selects a fermion formulation and
implements the Dirac operator, (b)~defines and constructs the
cross-projected propagator $\Scross=\Pew_{+}S\,\Pew_{-}$ on the
gauge-fixed backgrounds of \Cref{sec:gaugefix,sec:gribov}, and
(c)~verifies in the weak-coupling limit that $\Scross\ne0$ for a doublet
(D28) and $\Scross=0$ for a singlet (D29).

\begin{quote}\itshape
L/R caveat (second occurrence): The embedding does not, by itself,
distinguish $\SU(2)_L$ from $\SU(2)_R$; this requires external fermionic
input.
\end{quote}

\noindent
SL1 is \emph{addressed} here through the doublet/singlet structure, but
the caveat is not retired: the full distinction requires the
non-perturbative test of Paper 7.

\subsection{Fermion Formulation Selection (P6-F10)}
\label{ssec:formulation}

The Nielsen--Ninomiya theorem~\cite{NN1,NN2} forbids a local, Hermitian,
doubler-free, chirally symmetric lattice Dirac operator, so any
chirally-sensitive prediction requires a bypass. Two are standard. Domain
wall fermions (DWF)~\cite{DWF1,DWF2} introduce a fifth dimension of
extent $L_s$; chiral modes localise on opposite walls and exact chirality
is recovered as $L_s\to\infty$, with corrections $O(e^{-\alpha L_s})$.
Overlap fermions~\cite{NEU} satisfy the Ginsparg--Wilson (GW)
relation~\cite{GW} exactly at finite $a$ via
$D_N=(1/a)(1+\gfive\operatorname{sgn}(H_w))$, at higher cost. We adopt DWF
for \Cref{sec:fermionic} and Paper 7: the weak-coupling test needs only
approximate chirality (at $\be=20$, $L_s=8$, the residual mass
$m_{\mathrm{res}}\ll g^2$), and the DWF code carries forward to Paper 7.
The singlet vanishing (Lemma 5) is formulation-independent. Locked
parameters: $M_5=1.8$, $r=1$, $L_s=8$, $m_f=0.01$.

\subsection{Dirac Operator and the Ginsparg--Wilson Relation (P6-F11)}
\label{ssec:dirac}

The 4D Wilson kernel with domain wall height $M_5$ is
\begin{multline}
  D_W(M_5)_{xy,\alpha\beta,ij}=(4r+M_5)\delta_{xy}\delta_{\alpha\beta}\delta_{ij}\\
  -\tfrac12\sum_\mu\bigl[(r-\gamma_\mu)_{\alpha\beta}U^{ij}_\mu(x)\delta_{y,x+\hat\mu}
  +(r+\gamma_\mu)_{\alpha\beta}U^{\dagger ij}_\mu(x-\hat\mu)\delta_{y,x-\hat\mu}\bigr],
  \label{eq:wilsonkernel}
\end{multline}
where $U^{ij}_\mu(x)$ are the gauge-fixed links acting on the isospin
indices $i,j$ for a doublet, and replaced by $\delta_{ij}$ (no gauge
coupling) for a singlet. The 5D DWF operator
$D_{\mathrm{DWF}}[s,s']$ adds bulk hops $-\PR\delta_{s',s-1}$,
$-\PL\delta_{s',s+1}$ and boundary mass terms
$m_f\PL\delta_{s,0}\delta_{s',L_s-1}$,
$m_f\PR\delta_{s,L_s-1}\delta_{s',0}$, with
$\PL{}_{/R}=(1\mp\gfive)/2$. The physical propagator is the
boundary-to-boundary element
$S(x,y)=[D_{\mathrm{DWF}}^{-1}]_{s=0,s'=L_s-1}(x,y)$. At finite $L_s$,
\begin{equation}
  \gfive D_{\mathrm{eff}}+D_{\mathrm{eff}}\gfive
  =a\,D_{\mathrm{eff}}\gfive D_{\mathrm{eff}}+O(e^{-\alpha L_s}),
  \label{eq:GW}
\end{equation}
the GW relation, exact as $L_s\to\infty$. At $L_s=8$, $\be=20$, the
violation is negligible relative to the $O(g^2)$ signal of D28.

\subsection{Nielsen--Ninomiya Bypass}
\label{ssec:NN}

The Nielsen--Ninomiya theorem~\cite{NN1} states that correct continuum
limit, absence of doublers, and exact chiral symmetry cannot hold
simultaneously. DWF sacrifices exact chirality at finite $L_s$
($\{\gfive,D_{\mathrm{DWF}}\}\ne0$) and recovers it as
$O(e^{-\alpha L_s})$, a controlled bypass. SL2 is thereby partially
addressed: the weak-coupling test is an NN-bypass result, while full
chiral claims at non-perturbative coupling require $L_s\to\infty$ control
(Paper 7). Lemma 5 below is NN-independent.

\subsection{CG Solver (P6-F12)}
\label{ssec:cg}

$S(x,y)$ is obtained by solving $D_{\mathrm{DWF}}\Psi=\eta$ via conjugate
gradient on the normal equations $D^\dagger D\Psi=D^\dagger\eta$, with
relative residual tolerance $\varepsilon_{\mathrm{CG}}=10^{-10}$ and
even-odd preconditioning. For D28/D29 the source is at
$x_0=(0,0,0,0)$, wall $s=0$, Dirac index $\alpha=0$; the doublet source
sets isospin $i=0$ and is gauge-coupled, while the singlet source carries
no isospin index and propagates freely.

\subsection{The Cross-Projected Propagator (P6-F13)}
\label{ssec:scross}

The electroweak projectors (Paper 5 \S7.3) are
$\Pew_{+}=\operatorname{diag}(1,0)\otimes\Id_4$ and
$\Pew_{-}=\operatorname{diag}(0,1)\otimes\Id_4$, with
$\Pew_{+}+\Pew_{-}=\Id_8$, $(\Pew_{\pm})^2=\Pew_{\pm}$,
$\Pew_{+}\Pew_{-}=0$.

\begin{definition}[Cross-Projected Fermion Propagator]
For a doublet $\psi=(\psi_u,\psi_d)^{\mathsf T}$ with propagator $S(x,y)$,
\begin{equation}
  \Scross(x,y):=\Pew_{+}\,S(x,y)\,\Pew_{-}=S_{ud}(x,y),
  \label{eq:scrossdef}
\end{equation}
the upper-to-lower isospin block --- the channel sensitive to $W^{+}$
exchange at $O(g)$. The scalar norm is
$|\Scross|^2(x,y)=\sum_{\alpha\beta}|[\Scross(x,y)]_{\alpha\beta}|^2$.
\end{definition}

\begin{note}
The citation here is to Paper 5 \S7.3; the outline reference to
``Paper 3'' is a bibliography-key discrepancy.
\end{note}

\subsection{Lemma 5: Singlet Vanishing}
\label{ssec:lemma5}

\setcounter{lemma}{4}
\begin{lemma}[Singlet Vanishing of $\Scross$; paper-local]
For a fermion field $\psi$ transforming as an $\SU(2)$ singlet,
$\Scross(x,y)=\Pew_{+}S(x,y)\Pew_{-}=0$ identically, for all $x,y$, at
any $\be$, any lattice spacing, and any Dirac operator.
\end{lemma}

\begin{proof}
A singlet carries no $\SU(2)$ index and does not couple to $U_\mu$, so
its propagator is scalar in isospin space,
$S(x,y)=s(x,y)\,\Id_2\otimes M_{\mathrm{Dirac}}$, isospin-blind (the same
propagator in both isospin components). Since $\Pew_{\pm}$ act on the
$2$-dimensional isospin factor,
\[
  \Pew_{+}\bigl[s(x,y)\Id_2\otimes M\bigr]\Pew_{-}
  =s(x,y)\,[\Pew_{+}\Pew_{-}]\otimes M=0,
\]
using $\Pew_{+}\Pew_{-}=0$.
\end{proof}

\noindent
The proof uses neither chirality nor the form of the Dirac operator; the
result is universal. A non-zero singlet $\Scross$ is therefore a
definitive code diagnostic, motivating the inversion of D29.

\subsection{Weak-Coupling Argument: Doublet Non-Vanishing}
\label{ssec:doublet}

At $O(g)$ the $W^{+}$ field couples via $b=(U)_{01}=(ig/\sqrt2)W^{+}_\mu$,
so the doublet propagator acquires an off-diagonal block
$S_{ud}(x,x_0)=\sum_{z,\mu}S_0^{uu}(x,z)[gW^{+}_\mu(z)\Gamma_\mu]
S_0^{dd}(z,x_0)+O(g^2)$. Squaring and averaging,
$\langle|\Scross|^2\rangle\sim g^2\,V\times(\text{free convolution})>0$,
strictly positive. At $\be=20$, $g^2=0.2$ and the estimate is
$\langle|\Scross|^2\rangle\sim O(10^2)$, far above the D28 threshold. The
argument is weak-coupling only; the non-perturbative regime is Paper 7.

\subsection{Checks D28 and D29}
\label{ssec:D28D29}

\textbf{D28 (standard).} For a doublet source at $\be=20$, $N=6$, over
$\Ncfg\ge10$ gauge-fixed configurations,
$(1/\Ncfg)\sum_{\mathrm{cfg}}\sum_{x\ne x_0}|\Scross(x,x_0)|^2
>\varepsilon_{\mathrm{ferm}}=10^{-4}$. PASS confirms the perturbative
prediction (expected $O(10^2)$) --- the first numerical evidence that the
$b$-channel gateway (Proposition 5, Paper 5) is active in the fermionic
sector.

\textbf{D29 (inverted).} For a singlet source on the same configurations,
$(1/\Ncfg)\sum_{\mathrm{cfg}}\sum_{x\ne x_0}|\Scross|^2<10^{-10}$. PASS
(inverted) confirms Lemma 5. Per QS5, a non-zero result is not new
physics but a code error (singlet constructed with isospin coupling,
projector misapplied, or indices swapped); the tolerance $10^{-10}$ sits
four orders above the double-precision floor. The composite $\S8$
fermionic check passes only if both D28 and D29 pass.

\subsection{Scope Limit Status After \texorpdfstring{\Cref{sec:fermionic}}{Section 8}}
\label{ssec:scopestatus}

SL1 is addressed (fermionic input introduced; caveat continues through
\Cref{sec:discussion,sec:conclusions} pending Paper 7). SL2 is partially
addressed (DWF bypass at $L_s=8$; full chiral claims deferred). SL3 and
SL4 are unchanged; SL5 remains discharged.

%=======================================================================
\section{Discussion}
\label{sec:discussion}

Paper 6 advances the programme from a formal identification to a measured
propagator. We summarise the contribution through five conclusions and
update the open questions of Paper 5.

\paragraph{Conclusions.}
\textbf{(C1)} The gauge-fixing programme (P6-F1--P6-F9) and the
$\Gcrossgf$ measurement are specified, and Theorem 2 is established,
promoting the Paper 5 \S8.3 formal identification to an error-budgeted
quantity (numerical execution of checks D19--D26 pending).
\textbf{(C2)} Gribov copy variation is quantified (Proposition 7),
providing the first systematic uncertainty budget for the $b$-channel
propagator. \textbf{(C3)} $\DW(t)$ is extracted and a comparison to the $\SU(2)$
data of \cite{CM} is set up (\cite{Bo} enters only as a qualitative
$\SU(3)$ decoupling cross-reference); Part A (consistency, $\be\ge5$) is runnable and
Part B (direct comparison, $\be\le3.5$) awaits Scope D. \textbf{(C4)} The
sign structure of $\langle\Gcrossgf\rangle$ across $\be$ is resolved
(Corollary 6, \Cref{sec:signstructure}), answering OQ2. \textbf{(C5)} The
fermionic sector is set up: $\Scross$ is defined and the weak-coupling
doublet/singlet prediction tested (D28/D29), with the full
non-perturbative test deferred.

\paragraph{Updated open questions.}
OQ1 (Landau-gauge measurement) is addressed
(\Cref{sec:gaugefix}--\Cref{sec:comparison}); Gribov--Zwanziger region
analysis remains open. OQ2 (sign-definiteness) is characterised
(\Cref{sec:signstructure}); the continuum-limit extrapolation remains
open. OQ3 ($\SU(3)$ extension) remains open but is now more tractable
with the gauge-fixing infrastructure in place. OQ4 (chirality/fermion
content): the weak-coupling test is completed (\Cref{sec:fermionic}); the
strong-coupling chirality test is Paper 7's primary target. OQ5--OQ7
(electroweak mixing, higher-spin representations, full spacetime) are
unchanged and longer-range, save that the temporal links are now live
(OQ7 partially advanced).

%=======================================================================
\section{Conclusions}
\label{sec:conclusions}

\textbf{Result 1 (Theorem 2).} The Landau-gauge $W^{\pm}$ propagator
$\DW(t)$ is obtained from $\Gcrossgf$ via Theorem 2, promoting the Paper 5
\S8.3 formal identification to a fully specified, error-budgeted result
(numerical evaluation pending the verification programme) with an explicit
three-term error budget.

\textbf{Result 2 (Proposition 7 and Corollary 6).} Gribov copy variation
is quantified, and the sign structure of $\langle\Gcrossgf\rangle$ is
characterised as a function of $\be$ and volume, with a threshold $\bc$
and a $1/N^2$ extrapolation.

\textbf{Result 3 (Lemma 5; D28/D29).} The cross-projected propagator
$\Scross$ is defined and constructed; it vanishes identically for an
$\SU(2)$ singlet (Lemma 5), and the weak-coupling doublet prediction of
Paper 5 \S7.3 is verified perturbatively.

\textbf{NOT statement.} Paper 6 does not resolve the full Gribov problem,
does not perform the non-perturbative fermionic test, does not extend to
$\SU(3)$, does not introduce electroweak mixing or $U(1)_Y$, and does not
establish an $\SU(2)_L/\SU(2)_R$ distinction. The embedding does not, by
itself, distinguish $\SU(2)_L$ from $\SU(2)_R$; this requires external
fermionic input.

%=======================================================================
\appendix

\section{Updated Check Programme (D19--D29)}
\label{app:checks}

\renewcommand{\arraystretch}{1.2}
\begin{center}
\begin{tabular}{llp{0.38\textwidth}ll}
\toprule
ID & \S & Description & Type & Status\\
\midrule
D19 & 2 & $\FL$ convergence $\Theta<10^{-14}$ & standard & pending exec.\\
D20 & 2 & per-site Landau condition & standard & pending exec.\\
D21 & 2 & gauge-fixed link traces (3 parts) & standard & pending exec.\\
D22 & 3 & Gribov non-uniqueness $\Delta\Gcross>10^{-6}$ & inverted & pending exec.\\
D23 & 3 & best-copy stability under re-fixing & standard & pending exec.\\
D24 & 4 & $\Gcrossgf$ reproducibility $<10^{-12}$ & standard & pending exec.\\
D25 & 4 & weak-coupling consistency (4 parts) & standard & spec.; $\be{=}20$ pending\\
D26 & 5 & statistical convergence, $\mathrm{CV}<5\%$ & standard & pending exec.\\
D27a & 7 & sign positivity, weak coupling (Type-P) & standard & pending exec.\\
D27b & 7 & sign violation, strong coupling & inverted & deferred (Scope D)\\
D28 & 8 & doublet $|\Scross|^2>10^{-4}$ & standard & pending exec.\\
D29 & 8 & singlet $|\Scross|^2<10^{-10}$ & inverted & pending exec.\\
\bottomrule
\end{tabular}
\end{center}

\noindent
\textbf{Status key.} No D19--D29 result is executed in this draft: ``pending
exec.'' denotes a check that is fully specified and runnable on the relevant
available ensembles but whose numerical outcome is not yet reported here;
``deferred (Scope D)'' denotes a check awaiting ensemble generation; ``$\be{=}20$
pending'' denotes a check whose perturbative component is conditional on the
Group~E2 extended ensemble (\Cref{sec:ensembles}). Headline tables and
verification summaries are populated when these are run.

\noindent
Inverted checks (D22, D27b, D29) pass when the ``unexpected'' direction is
observed; each carries an explicit justification of what a standard-direction
result would mean physically (QS5). Any slope or sign anomaly is documented
with a resolution before a check is closed. Checks D1--D18 (Papers 4--5)
remain closed.

\section{Extended Ensemble and Gauge-Fixed Configuration Parameters}
\label{app:ensembles}

Table B.1 records the reprocessed Paper 4/5 ensembles (now gauge-fixed):
$\be$, $N$, $\Ncfg$, gauge-fixing algorithm, $\epsg$, Gribov copies per
configuration, and the best-copy $\FL$ value. Table B.2 records the new
extended ensembles ($\be\in\{2.0,2.5,3.0,3.5,20.0\}$, $N\in\{8,12,16\}$,
$\Ncfg\ge100$): thermalisation sweeps, separation, and acceptance.
Table B.3 records the fermionic parameters: formulation (DWF), domain wall
height $M_5=1.8$, Wilson parameter $r=1$, $L_s=8$, bare mass $m_f=0.01$,
$\varepsilon_{\mathrm{CG}}=10^{-10}$, and number of source points.
Table B.4 records the script assignments for D19--D29. Software versions:
SciPy 1.16.3, NumPy, and the FFT and sparse-linear-algebra dependencies
introduced for Fourier-accelerated fixing and the Dirac solver.

\paragraph{Paper 7 entry point.}
The full non-perturbative fermionic test at strong coupling is Paper 7;
the first available identifiers there are D30, Proposition 8, Corollary 7,
Theorem 3, and Lemma 6.

%=======================================================================
\begin{thebibliography}{99}

\bibitem{P4} T. Firestone, \emph{Paper 4 of the
$\mathbb{R}\otimes\mathbb{C}\otimes\mathbb{H}\otimes\mathbb{O}$ series}:
Cross-channel correlators in the $\Cl(0,2)$ embedding of lattice
$\SU(2)$.

\bibitem{P5} T. Firestone, \emph{Paper 5 of the
$\mathbb{R}\otimes\mathbb{C}\otimes\mathbb{H}\otimes\mathbb{O}$ series}:
The Two-Observable Decomposition in $\Cl(0,2)\cong\mathbb{H}$;
\S7.3 (conditional prediction) and \S8.3 (formal identification).

\bibitem{CM} A.~Cucchieri and T.~Mendes, \emph{Constraints on the
infrared behavior of the gluon propagator in Yang--Mills theories},
Phys. Rev. Lett. \textbf{100} (2008) 241601.

\bibitem{Bo} I.~L.~Bogolubsky, E.-M.~Ilgenfritz, M.~M\"uller-Preussker,
and A.~Sternbeck, \emph{Lattice gluodynamics computation of Landau-gauge
Green's functions in the deep infrared}, Phys. Lett. B \textbf{676}
(2009) 69--73.

\bibitem{MO} J.~E.~Mandula and M.~Ogilvie, \emph{Efficient gauge fixing
via overrelaxation}, Phys. Lett. B \textbf{248} (1990) 156--158.

\bibitem{Gri} V.~N.~Gribov, \emph{Quantization of non-Abelian gauge
theories}, Nucl. Phys. B \textbf{139} (1978) 1--19.

\bibitem{Zwa} D.~Zwanziger, \emph{Local and renormalizable action from
the Gribov horizon}, Nucl. Phys. B \textbf{323} (1989) 513--544.

\bibitem{Singer} I.~M.~Singer, \emph{Some remarks on the Gribov
ambiguity}, Commun. Math. Phys. \textbf{60} (1978) 7--12.

\bibitem{NN1} H.~B.~Nielsen and M.~Ninomiya, \emph{Absence of neutrinos
on a lattice. I. Proof by homotopy theory}, Nucl. Phys. B \textbf{185}
(1981) 20--40.

\bibitem{NN2} H.~B.~Nielsen and M.~Ninomiya, \emph{Absence of neutrinos
on a lattice. II. Intuitive topological proof}, Nucl. Phys. B
\textbf{193} (1981) 173--194.

\bibitem{DWF1} D.~B.~Kaplan, \emph{A method for simulating chiral
fermions on the lattice}, Phys. Lett. B \textbf{288} (1992) 342--347.

\bibitem{DWF2} Y.~Shamir, \emph{Chiral fermions from lattice boundaries},
Nucl. Phys. B \textbf{406} (1993) 90--106.

\bibitem{GW} P.~H.~Ginsparg and K.~G.~Wilson, \emph{A remnant of chiral
symmetry on the lattice}, Phys. Rev. D \textbf{25} (1982) 2649.

\bibitem{NEU} H.~Neuberger, \emph{Exactly massless quarks on the
lattice}, Phys. Lett. B \textbf{417} (1998) 141--144.

\end{thebibliography}

\end{document}
